\documentclass[answers, 10pt]{exam}
\usepackage[T1]{fontenc}
\usepackage{geometry}
\geometry{
	b5paper,
	total={136mm,210mm},
	left=20mm,
	top=20mm,
}

\usepackage[english]{babel}
\usepackage{amsmath,amssymb,amsthm}
\usepackage{mathrsfs,bm}
\usepackage{mathtools}
\usepackage{yhmath}
\usepackage{esint}

\usepackage{graphicx,epstopdf,caption,ulem}
\graphicspath{{figures/}}

\usepackage{float,subcaption,setspace,booktabs,multirow,supertabular,lscape,threeparttable,diagbox,extarrows}
\usepackage{xcolor}
\usepackage{colortbl}
\usepackage{multirow}
\usepackage{setspace}
\usepackage{ctex}
%\usepackage{CJK}
\usepackage{inputenc}
\usepackage{xeCJK}


\renewcommand{\solutiontitle}{\noindent\textbf{解:}\enspace}

\newcommand{\studentname}{请输入你的名字}
\newcommand{\studentid}{请输入你的学号}

\begin{document}

\pagestyle{headandfoot}
\runningheadrule
\firstpageheader{}{}{}
\runningheader{\small 高等数学（一）作业 \;\studentid\;\studentname}
{}
{\small 第\thepage 页}
\firstpagefooter{}{}{}
\runningfooter{}{}{}

\begin{center}
	{\Large\bf 高等数学（一）在线作业答题册}
	
	\medskip
	姓名：\studentname \qquad 学号：\studentid
\end{center}

\section*{解题格式范例}

\noindent {\bf 习题格式说明：}问题栏将标记 (选做) 与 (必做) 题, 并预留题号. 每章要求完成所有必做题, 课本与补充选做题{\color{red}各版块均达到半数以上}. 课本原题可自行抄写.

\smallskip
\noindent {\bf 解答文本要求：}证明陈述部分要求写明{\color{red}具体论证原理}. 计算推导部分要求{\color{red}写明每一步的推导理由}. 

\smallskip
\noindent {\bf 作业评分规则：}作业评分不计算正确率, {\color{red}仅考虑完成态度}. {\color{red}单纯复制黏贴解答}, {\color{red}书写逻辑不通或抄袭痕迹明显}等现象将会酌情扣分.

\smallskip
\noindent {\bf 格式规范模版：}

\begin{questions}
	
	\question {\bf 习题 1-2}{题 4.} (必做)
	
	(1) 若 $\displaystyle \lim_{n\to\infty} x_n=a$, 证明: $\displaystyle \lim_{n\to\infty} |x_n|=|a|$, 并举例说明反之不一定成立.
	
	(2) 若 $\displaystyle \lim_{n\to\infty} |x_n|=0$, 证明: $\displaystyle \lim_{n\to\infty} x_n=0$.
	
	\begin{solution}
		(1) 由定义 $\displaystyle \lim_{n\to\infty} x_n=a$, 则 $\forall \varepsilon>0$, $\exists N\in\mathbb{N}$, $\forall n\in\mathbb{N}$, $n>N$ 使得
		$$
		|x_n-a|<\varepsilon,
		$$
		那么对以上同样的 $\varepsilon>0$, $N\in\mathbb{N}$ 由三角不等式 $||\alpha|-|\beta||\le|\alpha-\beta|$ 必有 $\forall n\in\mathbb{N}$, $n>N$ 使得
		$$
		||x_n|-|a||\le|x_n-a|<\varepsilon,
		$$
		因此由义定可知 $\displaystyle \lim_{n\to\infty} |x_n|=|a|$ 得证.
		
		\smallskip
		举如下反例, 设: $x_n=(-1)^n$, 则显然 $\displaystyle \lim_{n\to\infty} |x_n|=1$ 而 $\displaystyle \lim_{n\to\infty} x_n$ 不存在.
		
		(2) 由定义 $\displaystyle \lim_{n\to\infty} |x_n|=0$, 则 $\forall \varepsilon>0$, $\exists N\in\mathbb{N}$, $\forall n\in\mathbb{N}$, $n>N$ 使得
		$$
		||x_n|-0|<\varepsilon \Rightarrow |x_n-0|<\varepsilon,
		$$
		同样由定义可知 $\displaystyle \lim_{n\to\infty} x_n=0$ 得证.
	\end{solution}
	
	\question {\bf 第一组补充题} (必做)
	
	运用重要极限
	$
	\displaystyle \lim_{n\to\infty} \left(1+\dfrac1{n}\right)^n = e
	$
	计算下列极限: \vspace{-0.5ex}
	$$
	\begin{aligned}
		& (1) \lim_{n\to\infty} \left(1-\frac1{n}\right)^n;
		&&(2) \ldots;
	\end{aligned} \vspace{-5ex}
	$$	
	\begin{solution}
		(1) 注意到({\it 此处运用平方差公式})
		$$
		\left(1-\frac1{n}\right) = \frac{\left(1+\frac1{n^2}\right)}{\left(1+\frac1{n}\right)},
		$$
		其中分母为重要极限
		$
		\displaystyle \lim_{n\to\infty} \left(1+\dfrac1{n}\right)^n = e.
		$
		分子又有
		$$
		\left(1+\frac1{n^2}\right)^n = \left(1+\frac1{n^2}\right)^{n^2\cdot\frac1{n}},
		$$
		注意到 $\left\{\left(1+\dfrac1{n^2}\right)^{n^2}\right\}$ 是 $\left\{\left(1+\dfrac1{n}\right)^{n}\right\}$ 的子列, 且后者为单调有界数列, 故
		$$
		1^n\le\left(1+\frac1{n^2}\right)^n\le e^{\frac1{n}},
		$$
		{\it 由夹逼原理}以及
		$
		\displaystyle \lim_{n\to\infty} \sqrt[n]{e} = 1
		$
		得
		$$
		\lim_{n\to\infty} \left(1+\dfrac1{n^2}\right)^n = 1.
		$$
		{\it 再由极限的四则运算法则}
		$$
		\lim_{n\to\infty} \left(1-\frac1{n}\right)^n = \frac{\displaystyle \lim_{n\to\infty} \left(1+\frac1{n^2}\right)^n}{\displaystyle \lim_{n\to\infty} \left(1+\frac1{n}\right)^n} = \frac1{e}.
		$$		
	\end{solution}
	
\end{questions}

\newpage

\section{函数}

\begin{questions}
	
	\question {\bf 第一组补充题} (必做)
	
	证明下列反三角函数等式或恒等式, 并解释其几何含义:
	
	\smallskip
	$
	\displaystyle
	(1)\; \arctan 1 + \arctan 2 + \arctan 3 = \pi; \qquad
	(2)\; \arctan x + \operatorname{arccot} x = \frac{\pi}{2};
	$
	
	\smallskip
	$
	\displaystyle
	(3)\; \arctan x + \arctan\left(\frac{1-x}{1+x}\right) = \frac{\pi}{4},\; (x>-1).
	$
	
	\begin{solution}
		在此处书写你的证明过程.
	\end{solution}
	
	\question {\bf 第一组补充题} (必做)
	
	证明下列反三角函数恒等式:
	
	\smallskip
	$
	\displaystyle
	(1)\; \arcsin x = \arctan\left(\frac{x}{\sqrt{1-x^2}}\right);		
	$
	
	\medskip
	$
	\displaystyle
	(2)\; 2\arcsin\sqrt{x} = \arccos(1-2x),\; x\in[0,1];
	$
	
	\medskip
	$
	\displaystyle
	(3)\; \arctan x \pm \arctan y = \arctan\left(\frac{x\pm y}{1\mp xy}\right).
	$
	
	\begin{solution}
		在此处书写你的证明过程.
	\end{solution}
	
	\question {\bf 第一组补充题} (必做)
	
	证明下列双曲函数恒等式:
	
	\smallskip
	$
	\displaystyle
	(1)\; \tanh(x\pm y) = \frac{\tanh x\pm\tanh y}{1\pm\tanh x\tanh y};
	$
	
	\medskip
	$
	\displaystyle
	(2)\; \cosh^4(2x)+4 = (4\sinh^4\!x+1)(4\cosh^4\!x+1).
	$
		
	\begin{solution}
		在此处书写你的证明过程.
	\end{solution}
	
	\question {\bf 第一组补充题} (必做)
	
	证明下列反双曲函数恒等式:
	
	\smallskip
	$
	\displaystyle
	(1)\; \operatorname{arsinh}x = \operatorname{artanh}\left(\frac{x}{\sqrt{1+x^2}}\right);
	$
	
	\medskip
	$
	\displaystyle
	(2)\; \ln|\tan x\,| = -\operatorname{artanh}(\cos2x);
	$
	
	\medskip
	$
	\displaystyle
	(3)\; \ln x = \operatorname{arsinh}\left(\frac{x^2-1}{2x}\right) = \operatorname{arcosh}\left(\frac{x^2+1}{2x}\right) = \operatorname{artanh}\left(\frac{x^2-1}{x^2+1}\right).
	$
	
	\begin{solution}
		在此处书写你的证明过程.
	\end{solution}
	
\end{questions}

\newpage

\section{数列极限}

\begin{questions}
	
	\question {\bf 习题 1-2}{题 3.} (选做)
	
	\begin{solution}
		在此处书写你的证明过程.
	\end{solution}
	
	\question {\bf 习题 1-2}{题 4.} (必做)
	
	\begin{solution}
		在此处书写你的证明过程.
	\end{solution}
	
	\question {\bf 习题 1-2}{题 5.} {\bf 习题 1-5}{题 4.} (必做)
	
	\begin{solution}
		在此处书写你的证明过程.
	\end{solution}
	
	\question {\bf 习题 1-2}{题 6.} (选做)
	
	\begin{solution}
		在此处书写你的证明过程.
	\end{solution}
	
	\question {\bf 习题 1-2}{题 7.} (选做)
	
	\begin{solution}
		在此处书写你的证明过程.
	\end{solution}
	
	\question {\bf 习题 1-6}{题 2.} (选做)
	
	\begin{solution}
		在此处书写你的证明过程.
	\end{solution}
	
	\question {\bf 习题 1-6}{题 3.} (选做)
	
	\begin{solution}
		在此处书写你的证明过程.
	\end{solution}
	
	\question {\bf 习题 1-6}{题 4.} (选做)
	
	\begin{solution}
		在此处书写你的证明过程.
	\end{solution}
	
\end{questions}

\begin{questions}
	
	\question {\bf 第一组补充题} (必做)
	
	运用极限定义解释以下命题:
	\begin{center}
		$\forall x\in\mathbb{R}$, $\exists\,\{a_n\mid a_n\in\mathbb{Q}\}$ 使得 $\displaystyle\lim_{n\to\infty} a_n = x$.
	\end{center}
	即任意实数均可通过一组有理点列逼近.
	
	\begin{solution}
		在此处书写你的证明过程.
	\end{solution}
	
	\question {\bf 第一组补充题} (必做)
	
	讨论满足下列条件时, 数列 $\{a_n+b_n\}$ 与 $\{a_nb_n\}$ 的敛散性:
	\begin{center}
		$(1)$ $\{a_n\}$与$\{b_n\}$中仅有一个收敛; \qquad $(2)$ $\{a_n\}$与$\{b_n\}$均发散.
	\end{center}
	
	\begin{solution}
		在此处书写你的证明过程.
	\end{solution}
	
	\question {\bf 第一组补充题} (必做)
	
	论述以下计算过程的错误, 并给出正确结果:
	$$
	\begin{aligned}
		\lim_{n\to\infty}\sum_{k=1}^{n} \frac1{n+k} = \lim_{n\to\infty} \frac1{n+1} + \lim_{n\to\infty} \frac1{n+2} + \cdots + \lim_{n\to\infty} \frac1{2n} = 0.
	\end{aligned} \vspace{-5ex}
	$$
	\begin{solution}
		在此处书写你的证明过程.
	\end{solution}
	
	\question {\bf 第一组补充题} (必做)
	
	计算下列极限: 
	$$
	\begin{aligned}
		& (1) \lim_{n\to\infty} \frac{n!}{n^n};
		&&(2) \lim_{n\to\infty} \frac1{\sqrt[n]{n!}};
		&&(3) \lim_{n\to\infty} (n+1)^{\frac1{\ln n}}.
	\end{aligned} \vspace{-5ex}
	$$
	\begin{solution}
		在此处书写你的证明过程.
	\end{solution}
	
	\question {\bf 第一组补充题} (必做)
	
	运用重要极限
	$
	\displaystyle \lim_{n\to\infty} \left(1+\dfrac1{n}\right)^n = e
	$
	计算下列极限: \vspace{-0.5ex}
	$$
	\begin{aligned}
		& (1) \lim_{n\to\infty} \left(1-\frac1{n}\right)^n;
		&&(2) \lim_{n\to\infty} \left(1+\frac{r}{n}\right)^n,\; r\in\mathbb{Q}; \\
		& (3) \lim_{n\to\infty} \left(1+\frac1{n}+\frac1{n^2}\right)^n;
		&&(4) \lim_{n\to\infty} n^2\left(\ln\left(1+\frac1{n}\right)-\frac1{n}\right).
	\end{aligned} 
	$$
	提示: 求解 $(4)$ 时, 注意到 $(\text{右侧数列是左侧数列的偶子列})$
	$$
	\frac1{2} \lim_{n\to\infty} n^2\left(\ln\left(1+\frac1{n}\right)-\frac1{n}\right) = \lim_{n\to\infty} 2n^2\left(\ln\left(1+\frac1{2n}\right)-\frac1{2n}\right)
	$$ \vspace{-5ex}
	\begin{solution}
		在此处书写你的证明过程.
	\end{solution}
	
	\question {\bf 第一组补充题} (必做)
	
	计算下列求和式或乘积式的极限:
	$$
	\begin{aligned}
		& (1) \lim_{n\to\infty} \sum_{k=0}^{n} r^k,\; (|r|<1);
		&&(2) \lim_{n\to\infty} \prod_{k=0}^{n} \left(1+r^{2^k}\right),\; (|r|<1); \\
		& (3) \lim_{n\to\infty} \prod_{k=2}^{n} \left(1-\frac{1}{k^2}\right);
		&&(4) \lim_{n\to\infty} \prod_{k=1}^{n} \left(1+\frac{k}{n^2}\right).
	\end{aligned}
	$$
	补充: 请尝试在不直接计算极限的情况下证明 $(1),(2)$ 两式相等.
	
	\medskip
	提示: 求解 $(4)$ 时, 先考察 $\dfrac{n^2}{k}\ln\left(1+\dfrac{k}{n^2}\right)$ 的极限.
	\begin{solution}
		在此处书写你的证明过程.
	\end{solution}
	
	\question {\bf 第一组补充题} (必做)
	
	证明: 数列 $\{\sin n\}$ 与 $\{\tan n\}$ 发散.
	
	\begin{solution}
		在此处书写你的证明过程.
	\end{solution}
	
\end{questions}

\begin{questions}
	
	\question {\bf 第二组补充题} (选做)
	
	运用 $\mathrm{Cauchy}$ 命题证明: 若 $\displaystyle\lim_{n\to\infty} a_n=L$ 则
	$$
	\lim_{n\to\infty} \sqrt[n]{a_1a_2\cdots a_n} = L.
	$$
	并计算
	$
	\displaystyle\lim_{n\to\infty} \frac{n}{\sqrt[n]{n!}}.
	$
	
	\begin{solution}
		在此处书写你的证明过程.
	\end{solution}
	
	\question {\bf 第二组补充题} (选做)
	
	证明: 若 $\displaystyle\lim_{n\to\infty} a_n=A$, $\displaystyle\lim_{n\to\infty} b_n=B$ 则
	$$
	\lim_{n\to\infty} \frac{a_1b_n+a_2b_{n-1}+\cdots+a_nb_1}{n} = AB.
	$$
	提示: 本题需要用到证明 $\mathrm{Cauchy}$ 命题时类似的技巧.
	
	\begin{solution}
		在此处书写你的证明过程.
	\end{solution}
	
	\question {\bf 第二组补充题} (选做)
	
	运用 $\mathrm{Stolz–Cesàro}$ 定理计算下列极限:
	$$
	\begin{aligned}
		& (1) \lim_{n\to\infty} \frac{n}{\sqrt[n]{n!}};
		&&(2) \lim_{n\to\infty} \frac{H_n}{\ln n},\; H_n\,\text{为调和数}; \\
		& (3) \lim_{n\to\infty} \frac1{\sqrt{n}}\left(1+\frac1{\sqrt2}+\cdots+\frac1{\sqrt{n}}\right);
		&&(4) \lim_{n\to\infty} \frac{1!+2!+\cdots+n!}{n!}.
	\end{aligned}
	$$  \vspace{-5ex}
	\begin{solution}
		在此处书写你的证明过程.
	\end{solution}
	
	\question {\bf 第二组补充题} (选做)
	
	计算下列极限:
	$$
	\begin{aligned}
		& (1) \lim_{n\to\infty} \frac{1^n+2^n+\cdots+n^n}{n^n}; \qquad
		&&(2) \lim_{n\to\infty} \frac1{n^2}\sum_{k=0}^{n}\ln\binom{n}{k}. 
	\end{aligned}
	$$
	提示: $(1)$ 不妨先考察 $\dfrac{(n-1)^n}{n^n}$, $\dfrac{(n-2)^n}{n^n},\,\cdots$ 的极限, 再尝试写出严谨的计算过程.
	
	\medskip
	$(2)$ 本题可连续两次运用 $\mathrm{Stolz–Cesàro}$ 定理求解.
	\begin{solution}
		在此处书写你的证明过程.
	\end{solution}
	
	\question {\bf 第二组补充题} (选做)
	
	证明: 数列 $\left\{\sqrt{1+\sqrt{2+\cdots+\sqrt{n}}}\,\right\}$ 收敛.
	
	\smallskip
	提示: 先证明 $\sqrt{n-1+\sqrt{n}}<1+\sqrt{n-1}$.
	
	\begin{solution}
		在此处书写你的证明过程.
	\end{solution}
	
	\question {\bf 第二组补充题} (选做)
	
	证明: 数列 $\sqrt7$, $\sqrt{7-\sqrt7}$, $\sqrt{7-\sqrt{7+\sqrt7}}$, $\sqrt{7-\sqrt{7+\sqrt{7-\sqrt7}}},\cdots$ 收敛, 并求其极限.
	
	\begin{solution}
		在此处书写你的证明过程.
	\end{solution}
	
	\question {\bf 第二组补充题} (选做) (Ramanujan 命题)
	
	证明: 数列 $\left\{\sqrt{1+2\sqrt{1+3\sqrt{1+\cdots+n}}}\,\right\}$ 收敛, 并求其极限.
	
	\smallskip
	提示: 本题为数学天才 $\mathrm{Ramanujan}$ 提出的经典问题之一, 建议查阅相关人物资料.
	
	\begin{solution}
		在此处书写你的证明过程.
	\end{solution}
	
	\question {\bf 第二组补充题} (选做)
	
	设数列首项 $a_0=1-a$ 且满足递推式
	$$
	a_{n+1}=\frac{a_n}{1-e^{-a_n}}-a,\;(0<a<1).
	$$
	证明: 数列 $\{a_n\}$ 收敛, 且其极限为方程 $f(x)=a$ 的实根, 其中
	$$
	f(x) = \frac{x}{e^x-1}.
	$$ \vspace{-5ex}
	
	\begin{solution}
		在此处书写你的证明过程.
	\end{solution}
	
	\question {\bf 第二组补充题} (选做) (算术--几何平均 Arithmetic–Geometric Mean)
	
	设对偶数列首项 $a_0=a$, $b_0=b$ 满足 $0<b<a$ 并由递推式
	$$
	a_n=\frac{a_{n-1}+b_{n-1}}{2},\quad b_n=\sqrt{a_{n-1}b_{n-1}}
	$$
	定义数列 $\{a_n\}$, $\{b_n\}$. 证明:
	$$
	\lim_{n\to\infty} a_n = \lim_{n\to\infty} b_n.
	$$
	此极限称为 $a,b$ 的{\bf 算术--几何平均}, 记为 $\operatorname{AGM}(a,b)$.
	
	\begin{solution}
		在此处书写你的证明过程.
	\end{solution}
	
\end{questions}
	
\end{document}